\subsection{Методология оценки качества}

Для объективной оценки разрабатываемого пайплайна и его сравнения с аналогами
необходимо создать методологию тестирования.
Процесс тестирования требует как подготовки валидационного датасета,
так и выбора подходящих метрик.

\subsubsection*{Реализация графа знаний в среде Obsidian и требования к данным}

Благодаря гибкости Obsidian и нативной поддержке ссылок,
создание ГЗ в нем будет выглядеть естественно.
В качестве сущностей будут выступать \texttt{Markdown}-файлы,
а в роли семантических связей~--- гиперссылки.
В данной работе будут также рассматриваться типизированные ссылки в Obsidian,
которые добавляет расширение Dataview.

Также необходимо установить ряд ограничений для сохранения каноничности графа,
а также упрощения работы с ним.
\begin{enumerate}
      \item Корпус текстов должен обладать высокой семантической связностью,
            то есть между файлами должны быть ярко выраженные семантические отношения
            (иерархические, причинно-следственные, функциональные и другие).
      \item Содержание каждого файла должно быть целостным и описательным.
      \item Атомарность сущностей: заметка должна представлять одну сущность.
      \item Пространство допустимых сущностей задается файлами в хранилище.
            При этом создание рефлексивных связей
            (когда заметка ссылается сама на себя) не допускается.
\end{enumerate}

Раскроем третье структурное ограничение подробнее.
Оно означает, что в содержании файла должна находиться информация только
об одной конкретной сущности, такой как личность, термин, событие, и так далее.
Содержание файла может включать в себя описание этой сущности, ее определение, историю,
факты, а также описание связей с другими сущностями.
Во-первых, такой подход делает хранилище максимально похожим на граф знаний.
Во-вторых, это убирает необходимость извлекать сущности,
так как файлы в хранилище уже формируют список допустимых сущностей.
В-третьих, данное ограничение соответствует концепции
атомарности заметок в Zettelkasten~\cite{Ahrens_2017},
где одна заметка содержит в себе ровно одну идею или концепцию.

\subsubsection*{Стратегия валидации}

Оценка такой комплексной системы от начала до конца
является некорректной, так как она может скрывать ошибки промежуточных этапов.
Например, если на этапе извлечения были получены нерелевантные чанки,
то этап типизации заведомо провалится, так как он не получит нужных данных.
В связи с этим валидация конвейера будет разделена на три этапа:
оценку извлечения контекста,
оценку типизации отношений,
а также оценку сгенерированного графа для использования в RAG-системах.
Такой подход даст возможность изолированного измерения качества
работы каждого модуля системы и его тонкой настройки,
что положительно скажется на итоговом качестве пайплайна.

\subsubsection*{Оценка извлечения контекста}

Оценку этапа извлечения (retrieval) можно сформулировать как задачу воспроизведения контекстов,
в которых упоминается сущность:
для каждой пары \((s, e)\), где \(s\)~--- файл-источник (\texttt{source\_file}),
а \(e\)~--- целевая сущность (\texttt{entity}),
система должна вернуть фрагмент из текста \(s\),
в достаточной мере совпадающий с эталонным контекстом из «золотого стандарта».
Извлеченный фрагмент должен полностью содержать эталонное предложение.

Важно заметить, что при больших чанках вероятность случайно захватить
нужное предложение выше.
Однако в реальных сценариях
извлекаемые фрагменты не могут быть слишком большими,
так как это приведет к размытию смыслов и потере деталей,
что негативно повлияет на качество извлечения.

Извлеченные фрагменты текста подразделяются на две категории:
\begin{itemize}
      \item Прямое (явное) упоминание.
            В данную группу входят все предложения,
            которые содержат полное или общепринятое краткое наименование.
            Например, предложение «Билл Гейтс соосновал Microsoft» явно содержит название компании.
      \item Кореферентное упоминание. Этот класс включает случаи,
            в которых прямое имя сущности заменено на
            другие языковые маркеры, требующие контекстного разрешения
            Например, в тексте «Microsoft была основана 4 апреля 1975.
            Билл Гейтс был одним из тех, кто ее соосновал» местоимение «ее» отсылает к Microsoft.
\end{itemize}

Для оценки необходимо создать специальный датасет, состоящий
из предложений с упоминаниями сущностей, списка упомянутых сущностей,
а также типа упоминания.

\subsubsection*{Оценка LLM-типизации}

Оценка типизации связей сводится к задаче классификации направленного
отношения для упорядоченной пары сущностей $(e_1, e_2)$.
При условии, что этап извлечения предоставил набор релевантных фрагментов,
большая языковая модель должна определить наличие связи от $e_1$ к $e_2$,
и присвоить ей один из представленных классов.
В случае присвоения недопустимого класса или
ответа модели о несвязанности сущностей
(классификации как отдельный класс \texttt{no\_relation}),
ссылка отбрасывается.

Для оценки качества классификации используются стандартные метрики~\cite{Manning_2008}:

\begin{itemize}
      \item Точность (Precision)~--- доля истинных связей среди
            всех предсказанных алгоритмом. Метрика важна для оценивания алгоритма на устойчивость к шуму.
            В случае создания множества ложноположительных связей граф засорится и потеряет свою ценность.
      \item Полнота (Recall)~--- доля истинных связей из датасета,
            которые алгоритм смог успешно идентифицировать и типизировать.
      \item Взвешенная F-мера (Weighted F\textsubscript{1}-score)~--- гармоническое
            среднее между точностью и полнотой. Эта метрика устойчива к дисбалансу в типах связей:
            зачастую в данных универсальные классы (такие как упоминания или дополнения)
            встречаются значительно чаще, чем, например, опровержение или противопоставление.
            Она будет использоваться в качестве основной оценки.
\end{itemize}

\subsubsection*{Полная оценка GraphRAG}

Для комплексной оценки ответов, сгенерированных LLM на основе контекста, предоставленного RAG,
используются специализированные фреймворки, такие как RAGAS~\cite{Es_2023}.
Их фундаментальной идеей является использование LLM как судей
для оценки ответов других языковых моделей,
что позволяет оценивать семантику ответов, а не использованные в них слова.

В данной работе будет использован фреймворк RAGAS.
Из метрик, предоставляемых им, будут подсчитываться:

\begin{itemize}
      \item Точность извлеченного контекста (Context Precision, CP).
            Данная метрика оценивает способность RAG систем
            извлекать релевантную информацию и ранжировать ее по мере соответствия.
            В RAGAS каждый извлеченный фрагмент текста оценивается
            на наличие и долю в нем нужной для ответа информации.
            Высокое значение CP говорит о предоставлении языковой модели
            релевантных данных с минимальной долей шума в них.
      \item Полнота извлеченного контекста (Context Recall, CR). Эта метрика оценивает RAG
            на умение собирать всю необходимую для ответа информацию.
            Для расчета этой метрики эталонный ответ
            разбивается на независимые атомарные утверждения, каждое из которых
            оценивается на возможность выведения его на основе извлеченных данных.
            Чем больше утверждений в эталонном ответе могут быть выведены,
            тем выше значение метрики.
      \item Достоверность ответа (Answer Faithfulness, F) направлена на измерение достоверности ответа,
            то есть степени,
            в которой сгенерированный ответ опирается исключительно на предоставленный контекст,
            а не на знания самой модели.
            Процесс оценки заключается в выделении всех тезисов в ответе и проверке
            возможности их подтверждения на основе предоставленных фрагментов текстов.
            Если модель использовала в ответе внешнюю информацию, то мера достоверности снижается.
      \item Корректность ответа (Answer Correctness, AC)
            измеряет семантическое и фактологическое соответствие
            сгенерированного ответа эталонному.
            Во фреймворке RAGAS она вычисляется как взвешенная сумма фактологического перекрытия
            (доля фактов, присутствующих как в ответе системы, так и в эталоне)
            и семантического сходства эмбеддингов двух ответов.
\end{itemize}

Использование такого набора метрик позволяет получить полную картину
работы пайплайна: от способности алгоритма GraphRAG предоставить релевантную информацию,
до способности модели использовать предоставленную информацию для генерации ответа.

\subsubsection*{Моделирование сложных сценариев для тестирования}

Чтобы проверить устойчивость алгоритма и его применимость в реальных условиях использования,
необходимо выявить наиболее сложные или важные ситуации,
с которыми встречаются алгоритмы построения ГЗ и GraphRAG.

\begin{itemize}
      \item Обнаружение сообществ. Сообщества~--- подграфы,
            вершины внутри которых плотно связаны между собой, но слабо с внешними.
            Они формируются из сущностей из одной тематики. Например,
            заметки с рецептами создают один кластер,
            а конспекты лекций по машинному обучению~--- другой.
            Современные алгоритмы ГЗ обычно автоматически
            их обнаруживают и создают для них отдельную вершину,
            в которой содержится краткая сводка~\cite{Edge_2024}.
            Эти вершины полезны при ответе на широкие вопросы,
            такие как
            «Какие темы из области математического анализа я изучал в этом месяце?»,
            а также для навигации алгоритмов поиска релевантных заметок.
            Они убирают нужду каждый запрос проходить по всем вершинам и
            обобщать их содержимое.
            Поскольку алгоритм, представленный в этой работе,
            не будет искать и создавать сообщества,
            уместно оценить возможность находить уже созданные.
            В рамках валидационного датасета вершины-сообщества будут созданы вручную.
      \item Проверка на устойчивость к шуму. В данном случае можно использовать омонимию,
            то есть слова или термины, совпадающие по написанию, но различные по смыслу.
            Такая структура заставит алгоритм использовать содержимое заметок для поиска и фильтрации.
            Также возможно создание графа из тесно связанных сущностей, что позволит
            протестировать способность алгоритма улавливать тонкие
            различия между типами связей в однородном тексте.
      \item Многошаговый поиск. Линейные цепочки из заметок, где
            каждая вытекает из предыдущей, проверяют
            эффективность поиска длинных логических путей.
\end{itemize}
